\chapter{Decentralized Multi-Agent Theorem Proving}\label{sec:decentralized}

The centralized MATP formulation presumes a benevolent planner with the authority to assign tasks to all agents. In practice, each agent controls its own compute, maintains private proof state, and decides independently what to prove and whether to share results. We now reformulate the problem as a multi-player game, first showing that decentralization introduces externalities that lead to inefficient equilibria, and then introducing a market mechanism that internalizes these externalities and aligns individual incentives with the social objective.


\subsection{Background: Multi-Agent Decision Making}\label{sec:posg-background}

We recall the standard frameworks for decentralized sequential decision-making under partial information.

\begin{definition}[Decentralized POMDP (Dec-POMDP)]
A discrete-time Dec-POMDP is a tuple
\[
\mathcal{D} = \Big(\mathcal{N},\;\mathsf{S},\;\{\mathcal{A}^\alpha\}_{\alpha\in\mathcal{N}},\;\{\mathsf{O}^\alpha\}_{\alpha\in\mathcal{N}},\;T,\;T^0,\;\{Z^\alpha\}_{\alpha\in\mathcal{N}},\;r\Big),
\]
where $\mathcal{N}$ is a finite set of agents, $\mathsf{S}$ is the state space, $\mathcal{A}^\alpha$ and $\mathsf{O}^\alpha$ are the action and observation spaces for agent $\alpha$, $T$ and $T^0$ are the transition and initial-state kernels, $Z^\alpha(\cdot\mid s,a,s')$ is agent $\alpha$'s observation kernel, and $r:\mathsf{S}\times\mathcal{A}\to\mathbb{R}$ is a \emph{shared} reward function (where $\mathcal{A}:=\prod_\alpha \mathcal{A}^\alpha$). The defining feature is that all agents share the same reward $r$, making the problem \emph{cooperative}: agents wish to maximize the same objective but must do so using only their private observation histories.
\end{definition}

\begin{definition}[Partially Observable Stochastic Game (POSG)]
A POSG generalizes the Dec-POMDP by replacing the shared reward with \emph{individual} reward functions $\{r^\alpha\}_{\alpha\in\mathcal{N}}$, where $r^\alpha:\mathsf{S}\times\mathcal{A}\to\mathbb{R}$. A Dec-POMDP is a POSG in which $r^\alpha = r^\beta$ for all $\alpha,\beta$. When the reward functions differ, agents have potentially conflicting objectives, and the appropriate solution concept shifts from cooperative optimality to game-theoretic equilibrium.
\end{definition}

Each agent $\alpha$ selects actions according to a policy $\pi^\alpha$ that maps its private action--observation history $h^\alpha_t := (o^\alpha_0,a^\alpha_0,\ldots,o^\alpha_t)$ to a distribution over $\mathcal{A}^\alpha$. A joint policy is $\pi = (\pi^\alpha)_{\alpha\in\mathcal{N}}$, and agent $\alpha$'s expected discounted return under $\pi$ is
\[
R^\alpha(\pi) := \mathbb{E}^{\pi}\!\left[\sum_{t=0}^{\infty}\gamma^t\, r^\alpha(s_t,a_t)\right].
\]

\begin{definition}[Nash equilibrium]
A joint policy $\pi^*$ is a \emph{Nash equilibrium} (NE) of a POSG if no agent can improve its expected return by unilaterally deviating:
\[
R^\alpha(\pi^{*\alpha},\pi^{*,-\alpha}) \;\ge\; R^\alpha(\tilde\pi^\alpha,\pi^{*,-\alpha})
\qquad\text{for all }\alpha\in\mathcal{N},\; \tilde\pi^\alpha\in\Pi^\alpha.
\]
\end{definition}

\begin{definition}[Social welfare and Price of Anarchy]
The \emph{social welfare} of a joint policy $\pi$ is $\mathrm{SW}(\pi) := \sum_{\alpha\in\mathcal{N}} R^\alpha(\pi)$. The \emph{Price of Anarchy} (PoA) is
\[
\mathrm{PoA} := \frac{\sup_\pi \mathrm{SW}(\pi)}{\inf_{\pi^*\in\mathrm{NE}} \mathrm{SW}(\pi^*)},
\]
measuring the worst-case efficiency loss from strategic behavior relative to the social optimum.
\end{definition}


\subsubsection{Complexity of decentralized control.}

\begin{proposition}[Complexity of Dec-POMDP and POSG]\label{prop:dec-complexity}
\leavevmode
\begin{enumerate}
    \item Computing an optimal joint policy for a finite-horizon Dec-POMDP is NEXP-complete (Bernstein, Givan, and Zilberstein, 2002).
    \item Computing a Nash equilibrium of a finite POSG is at least PPAD-hard (by reduction from Nash equilibrium computation in normal-form games).
\end{enumerate}
\end{proposition}

Thus, even the cooperative problem (Dec-POMDP) is vastly harder than the centralized problem (POMDP, which is PSPACE-complete). Decentralization incurs a fundamental computational cost.


\subsection{The Decentralized MATP Game}\label{sec:dec-matp}

We now formulate the \emph{Decentralized MATP} as a POSG by distributing the centralized planner's authority to the individual agents.

\begin{definition}[Decentralized MATP POSG]\label{def:dec-matp}
Given the same primitives as the centralized MATP (agent set $\mathcal{N}$, formula space $\mathcal{F}$, proof and conjecture kernels, library structure), the \emph{Decentralized MATP} is the POSG
\[
\mathcal{G}_{\mathrm{MATP}} = \Big(\mathcal{N},\;\mathsf{S},\;\{\mathcal{A}^\alpha\}_{\alpha\in\mathcal{N}},\;\{\mathsf{O}^\alpha\}_{\alpha\in\mathcal{N}},\;T,\;T^0,\;\{Z^\alpha\}_{\alpha\in\mathcal{N}},\;\{r^\alpha\}_{\alpha\in\mathcal{N}},\;\gamma\Big)
\]
where:
\begin{itemize}
    \item \textbf{State space $\mathsf{S}$} is identical to the centralized MATP (private libraries, public library, job descriptors, query models, cumulative effort).
    
    \item \textbf{Action spaces $\mathcal{A}^\alpha$} are identical to the per-agent action spaces of the centralized MATP: each agent independently selects $\mathrm{Prove}$, $\mathrm{Conj}$, $\mathrm{Pub}$, $\mathrm{Qry}$, or $\bot$.
    
    \item \textbf{Observations $\mathsf{O}^\alpha$} are agent-specific: agent $\alpha$ observes its own private library $\mathcal{L}^\alpha_t$, its own portfolio/balance, its own query responses $\Omega^\alpha_t$, and the public library $\mathcal{L}^0_t$. Crucially, agent $\alpha$ does \emph{not} observe other agents' private libraries, job statuses, or cumulative effort maps.
    
    \item \textbf{Transition kernel $T$} is identical to the centralized MATP (the physics of proving, conjecturing, and publishing are unchanged).
    
    \item \textbf{Individual rewards $r^\alpha$} encode each agent's private objective. We defer the precise specification of $r^\alpha$ to the discussion below, noting that the choice of reward structure is the central design variable.
\end{itemize}
\end{definition}

The key structural change from the centralized MATP is that there is no planner: each agent $\alpha$ selects its own action $a^\alpha_t$ based on its own observation history $h^\alpha_t$. The joint action $a_t = (a^\alpha_t)_{\alpha\in\mathcal{N}}$ emerges from the independent, simultaneous choices of all agents.


\subsubsection{Cooperative baseline: Dec-POMDP.}

As a baseline, one may set $r^\alpha = r$ for all $\alpha$, recovering a Dec-POMDP in which all agents cooperatively maximize the centralized MATP objective (e.g., minimizing the time to fully resolve $\Phi^*$). In this case, communication occurs solely through publication: an agent that publishes a resolved sequent effectively ``communicates'' its result to all others via the shared library. The Dec-POMDP formulation captures the difficulty of coordination under communication constraints, but assumes agents are willing to cooperate---they share the social objective and have no private interests.

\begin{remark}[Communication via publication]
In the Dec-POMDP, the only communication channel between agents is the shared library $\mathcal{L}^0_t$. An agent's publication action $(\mathrm{Pub},(C,R))$ is the sole mechanism by which private information enters the public domain. This is a severely constrained communication model: agents cannot send arbitrary messages, and publication is ``all-or-nothing'' (either a result is in the shared library or it is not). The Dec-POMDP's NEXP-completeness (Proposition~\ref{prop:dec-complexity}) reflects, in part, the difficulty of coordinating under such constraints.
\end{remark}


\subsubsection{Individual rewards and the externality problem.}

In practice, agents are not purely cooperative. Each agent has its own research agenda, computational resources, and incentives. A more realistic model assigns each agent an individual reward that depends on the agent's own contributions and on the state of the shared library.

Consider a natural class of \emph{credit-based} reward functions: agent $\alpha$'s reward is determined by the sequents that $\alpha$ has proven and that are credited to $\alpha$ in the metadata $\Theta^0_t$.

\begin{definition}[Credit-based reward]\label{def:credit-reward}
Fix a target valuation $v:\mathcal{S}\to\mathbb{R}_{\ge 0}$ assigning value to fully resolved sequents. Agent $\alpha$'s credit-based per-step reward is
\[
r^\alpha_{\mathrm{credit}}(s_t,a_t) := \sum_{\substack{s\in\mathcal{FRS}^0_{t+1}\setminus\mathcal{FRS}^0_t}} v(s)\cdot\mathbf{1}[\Theta^0_{t+1}(s) = (\cdot,\alpha)],
\]
i.e., agent $\alpha$ is rewarded for each newly fully resolved sequent in the shared library that is credited to $\alpha$.
\end{definition}

Under credit-based rewards, agents have incentive to prove valuable sequents and receive credit. However, this reward structure introduces a fundamental externality.


\subsection{The Publication Externality}\label{sec:externality}

We now show that the decentralized MATP with credit-based rewards exhibits a \emph{publication externality}: an agent's publication of a resolved sequent creates social value (by enabling other agents' proofs) but generates no direct reward for the publisher.

\begin{definition}[Publication externality]\label{def:pub-externality}
A \emph{publication externality} exists at state $s_t$ if there exists an agent $\alpha$ and a publication action $a^\alpha = (\mathrm{Pub},(C,R))$ such that:
\begin{enumerate}
    \item Publishing strictly increases some other agent $\beta$'s expected future reward:
    \[
    \mathbb{E}[R^\beta_{\ge t+1} \mid s_t,\; a^\alpha = (\mathrm{Pub},(C,R))] > \mathbb{E}[R^\beta_{\ge t+1} \mid s_t,\; a^\alpha = \bot],
    \]
    where $R^\beta_{\ge t+1} := \sum_{\tau=t+1}^\infty \gamma^{\tau-t-1} r^\beta(s_\tau,a_\tau)$.
    \item Publishing does not increase $\alpha$'s own expected future reward:
    \[
    \mathbb{E}[R^\alpha_{\ge t+1} \mid s_t,\; a^\alpha = (\mathrm{Pub},(C,R))] \le \mathbb{E}[R^\alpha_{\ge t+1} \mid s_t,\; a^\alpha = \bot].
    \]
\end{enumerate}
\end{definition}

\begin{theorem}[Existence of publication externalities]\label{thm:externality}
In any Decentralized MATP POSG with $|\mathcal{N}|\ge 2$, credit-based rewards, and non-trivial latent dependencies between agents' targets, publication externalities exist.
\end{theorem}

\begin{proof}
We construct a witnessing instance. Let $\mathcal{N} = \{\alpha,\beta\}$. Let the target sets be $\Phi^*_\alpha = \{s_\alpha\}$ and $\Phi^*_\beta = \{s_\beta\}$, with $v(s_\alpha) = v(s_\beta) = 1$ and $v(s) = 0$ otherwise. Suppose the latent dependency structure satisfies $\delta^*(s_\beta) = \{\phi_\alpha\}$, where $\phi_\alpha$ is the consequent of $s_\alpha$. Thus $s_\beta$ cannot be fully resolved until $s_\alpha$ is also resolved in the shared library.

At some time $t$, suppose $\alpha$ has privately resolved $s_\alpha$ (i.e., $s_\alpha\in\mathcal{RS}^\alpha_t$) but has not published it ($s_\alpha\notin\mathcal{RS}^0_t$). Agent $\alpha$'s own target $s_\alpha$ is already resolved in $\alpha$'s private library, so $\alpha$'s credit-based reward for $s_\alpha$ is earned only when $s_\alpha$ enters $\mathcal{FRS}^0_t$---which requires publication. Under credit-based rewards, $\alpha$ receives $v(s_\alpha) = 1$ upon publishing $s_\alpha$ (since $\delta^*(s_\alpha)=\varnothing$ makes it immediately fully resolved upon publication). So publishing actually does help $\alpha$ in this case.

To create a pure externality, modify the instance so that $\alpha$'s target is $s'_\alpha$ with $\delta^*(s'_\alpha) = \varnothing$ and $s'_\alpha \neq s_\alpha$, while $\alpha$ has also privately resolved $s_\alpha$ as a byproduct of working toward $s'_\alpha$. Now $\alpha$ receives credit for $s'_\alpha$ (its own target) but not for $s_\alpha$ (which is not in $\alpha$'s target set and has $v(s_\alpha) = 0$ from $\alpha$'s perspective). Publishing $s_\alpha$ costs $\alpha$ one timestep but provides $\alpha$ no reward (since $v(s_\alpha) = 0$ for $\alpha$). However, publishing $s_\alpha$ enables $\beta$ to fully resolve $s_\beta$ (since $s_\beta$'s dependency on $\phi_\alpha$ is satisfied). Thus:
\begin{itemize}
    \item $\alpha$'s reward from publishing $s_\alpha$: $0$ (no credit, cost of one timestep).
    \item $\beta$'s reward from $\alpha$ publishing $s_\alpha$: $v(s_\beta) = 1$ (enables full resolution of $\beta$'s target).
\end{itemize}
This satisfies Definition~\ref{def:pub-externality}.
\end{proof}

\begin{corollary}[Free riding and underinvestment]\label{cor:free-riding}
Under credit-based rewards, the unique Nash equilibrium of the externality instance in Theorem~\ref{thm:externality} has $\alpha$ not publishing $s_\alpha$, leading to $\beta$ failing to fully resolve $s_\beta$. The social welfare at this NE is strictly lower than the social optimum (in which $\alpha$ publishes $s_\alpha$, enabling $\beta$ to resolve $s_\beta$).
\end{corollary}

The publication externality is the core market failure in decentralized theorem proving: agents who produce useful intermediate results have no incentive to share them. This leads to duplication of effort (multiple agents independently proving the same lemma), underinvestment in ``infrastructure'' lemmas (results useful to many but credited to none), and withholding of results (agents hoarding private proofs to maintain competitive advantage).

\begin{remark}[Beyond credit-based rewards]
The publication externality is not an artifact of the credit-based reward structure. Any individual reward function that fails to compensate agents for the social value of their publications will exhibit a similar externality. The fundamental issue is that publication creates positive externalities---benefits to other agents that are not captured by the publisher's reward. Eliminating this externality requires a mechanism that transfers value from beneficiaries to publishers: a market.
\end{remark}


\subsection{Market-Mediated MATP}\label{sec:market-matp}

We now introduce a market mechanism that internalizes the publication externality by enabling agents to trade state-contingent contracts whose settlement is triggered by publicly verifiable events (sequent resolution in the shared library).

\subsubsection{Collateralized bilateral contract market.}

We equip the Decentralized MATP POSG with a \emph{bilateral contract market} operating as follows.

A \emph{contract type} is a triple $\chi = (s, T, \ell) \in \mathcal{CS}_0 \times \mathbb{N} \times [0,1]$, representing a normalized unit-notional claim on the event that sequent $s$ becomes publicly resolved by deadline $T$, with maximum loss parameter $\ell$. Each contract type induces two complementary positions, the \emph{long} side $\chi^+$ and the \emph{short} side $\chi^-$, with terminal payoffs:
\[
\Pi(\chi^+) = \begin{cases} 1-\ell & \text{if } s\in\mathcal{RS}^0_T, \\ -\ell & \text{otherwise,} \end{cases}
\qquad
\Pi(\chi^-) = -\Pi(\chi^+).
\]

Contracts are traded via \emph{public offers} $o = (\chi, \sigma, q, p)$, where $\sigma\in\{+,-\}$ is the side offered, $q\in\mathbb{N}^+$ is quantity, and $p\in[0,1]$ is the per-unit price. When agent $\beta$ accepts quantity $q'\le q$ from an offer posted by $\alpha$, the position $q'\chi^\sigma$ transfers to $\beta$'s holdings and cash $q'p$ transfers from $\beta$ to $\alpha$.

Contracts are always issued as matched pairs $(q\chi^+, q\chi^-)$, ensuring zero net creation of wealth. A \emph{collateral constraint} requires that every agent's portfolio satisfies $\underline{B}^\alpha \ge 0$ at all times, where $\underline{B}^\alpha := c^\alpha - \sum_{x\in H^\alpha\cup O^\alpha} \sup_\omega(-\Pi_\omega(x))_+$ is the minimum marked balance (cash minus worst-case liabilities). This ensures no agent can default.

At each timestep $t$, all contracts with deadline $T = t$ are \emph{settled}: payoffs are applied to holdings based on the public library state $\mathcal{RS}^0_t$, and settled contracts are removed.


\subsubsection{Market-augmented reward.}

In the market-mediated MATP, each agent $\alpha$'s reward is the change in portfolio balance:
\[
r^\alpha_{\mathrm{mkt}}(s_t,a_t) := B^\alpha_{t+1} - B^\alpha_t,
\]
where $B^\alpha_t = \beta(\mathcal{P}^\alpha_t)$ is the canonical balance of $\alpha$'s portfolio (cash plus mark-to-market value of positions). This reward naturally captures trading profits, proof costs, and contract settlement payoffs.


\begin{definition}[Market-Mediated MATP (VAMP)]\label{def:vamp-game}
The \emph{Verifiable Allocation Market POSG} (VAMP) is the POSG obtained by augmenting the Decentralized MATP with the bilateral contract market:
\[
\mathcal{G}_{\mathrm{VAMP}} = \Big(\mathcal{N},\;\mathcal{W},\;\{\mathcal{A}^\alpha\}_{\alpha\in\mathcal{N}},\;\{\mathsf{O}^\alpha\}_{\alpha\in\mathcal{N}},\;\mathcal{T},\;\{r^\alpha_{\mathrm{mkt}}\}_{\alpha\in\mathcal{N}},\;\gamma,\;b_0,\;\mathscr{M}\Big),
\]
where the world state $\mathcal{W}$ includes the market state $\mathcal{M}_t = (\{\mathcal{P}^\alpha_t\}_\alpha, \mathcal{ML}_t)$, agent actions include market actions (posting offers, accepting offers) alongside proof/conjecture/publication actions, and the transition kernel incorporates the two-stage construction (non-market transitions followed by market transitions) as specified by the market mechanism $\mathscr{M}$. The full formal definition is given in Chapter~\ref{chap:vamp}.
\end{definition}


\subsection{Internalization of Externalities}\label{sec:internalization}

We now show that the bilateral contract market internalizes the publication externality identified in Theorem~\ref{thm:externality}.

\begin{theorem}[Market internalization of publication externalities]\label{thm:internalization}
Consider the VAMP game $\mathcal{G}_{\mathrm{VAMP}}$ with the bilateral contract market. Suppose agent $\beta$ values the public resolution of sequent $s$ by time $T$ at $v_\beta > 0$ (in the sense that $\beta$'s expected future reward increases by $v_\beta$ if $s\in\mathcal{RS}^0_T$). Then:
\begin{enumerate}
    \item \textbf{Demand expression.} Agent $\beta$ can express demand for the resolution of $s$ by purchasing long contracts $\chi^+ = (s,T,\ell)^+$ at a price $p < v_\beta$. This is individually rational for $\beta$ whenever the market price is below $\beta$'s valuation.
    
    \item \textbf{Proof incentive.} If agent $\alpha$ can resolve $s$ at expected cost $c_\alpha$, then $\alpha$ can profitably buy long contracts at any price $p < 1 - \ell - c_\alpha$, then prove and publish $s$, collecting the settlement payoff $1-\ell$ per contract. The publication externality is internalized: $\alpha$'s reward from proving and publishing now includes the contract profit $(1-\ell) - p - c_\alpha > 0$.
    
    \item \textbf{Value transfer.} In any trade, the contract price $p$ transfers value from the buyer (who benefits from resolution) to the seller of the complementary side. At equilibrium, the long price $p^*$ satisfies
    \[
    c_{\min} \;\le\; 1 - \ell - p^* \;\le\; v_{\max},
    \]
    where $c_{\min} = \min_\alpha c_\alpha$ is the cheapest prover's cost and $v_{\max} = \max_\beta v_\beta$ is the highest valuation. Trades occur whenever there are gains from trade ($c_{\min} < v_{\max}$), and the market price reflects the marginal cost of proof provision.
\end{enumerate}
\end{theorem}

\begin{proof}[Proof sketch]
(1) Agent $\beta$ buys $\chi^+$ at price $p$. If $s$ is resolved by $T$, $\beta$ receives $1-\ell$ from settlement. $\beta$'s expected profit from the contract is $\Pr(s\in\mathcal{RS}^0_T)(1-\ell) - \ell\cdot\Pr(s\notin\mathcal{RS}^0_T) - p$. If $\beta$'s valuation $v_\beta$ exceeds $p$, $\beta$ has positive expected surplus from purchasing.

(2) Agent $\alpha$ buys $\chi^+$ at price $p$, then proves $s$ at cost $c_\alpha$ and publishes. Settlement yields $1-\ell$ per contract. $\alpha$'s net payoff is $(1-\ell) - p - c_\alpha$. This is positive whenever $p < 1-\ell-c_\alpha$. Crucially, $\alpha$ now has a direct financial incentive to both prove \emph{and} publish $s$: the settlement payoff is triggered by public resolution ($s\in\mathcal{RS}^0_T$), which requires publication.

(3) At equilibrium, the long price must make the marginal buyer indifferent ($p^* \approx v_{\mathrm{marginal}} - (1-\ell)$, adjusted for resolution probability) and the marginal prover indifferent ($p^* \approx (1-\ell) - c_{\mathrm{marginal}}$). Trades occur when these ranges overlap, i.e., when some prover's cost is below some buyer's valuation.
\end{proof}

\begin{remark}[Comparison to the externality instance]
Revisiting the instance from Theorem~\ref{thm:externality}: without a market, $\alpha$ has no incentive to publish $s_\alpha$ (zero reward, positive cost). With the market, $\beta$ can buy long contracts on $s_\alpha$, creating a financial incentive for any agent (including $\alpha$) to prove and publish $s_\alpha$. The market price reflects $\beta$'s valuation, and the resulting trade transfers value from $\beta$ (beneficiary) to the prover (provider), eliminating the externality.
\end{remark}

\begin{remark}[Collateralization and manipulation]\label{rem:collateral}
The collateral constraint $\underline{B}^\alpha \ge 0$ ensures that agents cannot take on unbounded liabilities, preventing default risk. This is crucial for incentive alignment: without collateralization, agents could sell short contracts on sequents they intend to prevent from being resolved (a form of market manipulation). Full collateralization bounds the worst-case loss for every agent, ensuring that all positions are credibly backed. However, collateralization also constrains participation: agents with limited capital cannot take large positions, potentially limiting market liquidity. The tradeoff between manipulation-resistance and liquidity is a standard theme in market design.
\end{remark}


\subsection{Equilibrium Analysis}\label{sec:equilibria}

We now analyze the game-theoretic properties of the VAMP game.

\subsubsection{Existence of equilibria.}

\begin{proposition}[Existence of Nash equilibria]\label{prop:ne-existence}
If the VAMP game $\mathcal{G}_{\mathrm{VAMP}}$ has finite state, action, and observation spaces, compact strategy sets, and continuous payoff functions (in the policy topology), then a Nash equilibrium exists in mixed (behavioral) strategies.
\end{proposition}

\begin{proof}[Proof sketch]
By Glicksberg's theorem, every finite game with compact strategy sets and continuous payoffs has a Nash equilibrium in mixed strategies. The VAMP's finite-horizon truncation satisfies these conditions. For the infinite-horizon discounted game, existence follows from the standard fixed-point argument on the space of behavioral strategies, applying Kakutani's theorem to the best-response correspondence.
\end{proof}

\begin{remark}[Markov perfect equilibria]
For computational and analytical tractability, one often restricts to \emph{Markov perfect equilibria} (MPE), in which each agent's policy depends on the current state only (not the full history). In the VAMP, a state-dependent policy $\pi^\alpha(s_t)$ conditions on agent $\alpha$'s current observation of the world state (private library, portfolio, public library). Under standard regularity conditions, MPE exist and are refinements of Nash equilibria.
\end{remark}


\subsubsection{Price of Anarchy.}

The Price of Anarchy measures the efficiency loss from strategic behavior relative to the social optimum. In our setting, the social optimum corresponds to the centralized MATP solution (a single planner maximizing social welfare), and the PoA compares the worst NE of the VAMP game against this benchmark.

\begin{proposition}[PoA lower bound]\label{prop:poa-lower}
The Price of Anarchy of the Decentralized MATP POSG \emph{without} a market (credit-based rewards) is unbounded: for any $M > 0$, there exists an instance with $\mathrm{PoA} > M$.
\end{proposition}

\begin{proof}
Consider the externality instance from Theorem~\ref{thm:externality} with $n$ agents, where each agent $\alpha_i$ privately resolves a sequent $s_i$ that is a dependency of the next agent $\alpha_{i+1}$'s target. Without a market, no agent has incentive to publish, so only agent $\alpha_1$'s target is achievable in the worst NE. The social optimum achieves all $n$ targets. Thus $\mathrm{PoA} \ge n$, which is unbounded as $n\to\infty$.
\end{proof}

\begin{proposition}[PoA with market]\label{prop:poa-market}
In the VAMP game with the bilateral contract market and sufficient initial liquidity, the Price of Anarchy is bounded. Specifically, if all agents have sufficient capital to participate in the market (i.e., initial balances $c^\alpha_0 \ge \bar c$ for some threshold $\bar c$ depending on the instance), then at any Nash equilibrium, every sequent $s$ with social value exceeding the cheapest resolution cost is resolved. In particular, the PoA is at most a constant depending on the market parameters $(\ell, \gamma)$.
\end{proposition}

\begin{proof}[Proof sketch]
At any NE, if a sequent $s$ has positive social surplus (some agent values its resolution above the cheapest prover's cost), then a profitable deviation exists for some agent: buy long, prove, publish, collect settlement. This deviation is unilateral and profitable, contradicting NE if $s$ remains unresolved. The bound on PoA follows from bounding the fraction of social value that is captured by market frictions (transaction costs implicit in the bid-ask spread, the loss parameter $\ell$, and discounting).
\end{proof}


\subsubsection{Computational complexity of equilibria.}

\begin{proposition}[Complexity of VAMP equilibria]\label{prop:vamp-complexity}
\leavevmode
\begin{enumerate}
    \item Computing an exact Nash equilibrium of the VAMP game is PPAD-hard (by reduction from bimatrix games, since the market introduces a general-sum structure even with two agents).
    \item Computing an optimal Dec-POMDP solution (the cooperative benchmark) is NEXP-complete.
    \item The VAMP's state space, augmented with market state (portfolios, order book), is exponentially larger than the centralized MATP state space. Exact equilibrium computation is therefore doubly intractable: both the game-theoretic and the state-space dimensions are hard.
\end{enumerate}
\end{proposition}


\subsection{Intractibility Consequences}\label{sec:game-summary}

% The analysis of this section establishes the following chain of results:

% \begin{enumerate}
%     \item The centralized MATP (POMDP/MDP) provides the social-optimum benchmark but assumes unrealistic centralized control.
%     \item The decentralized MATP (Dec-POMDP/POSG) removes centralized control, introducing both coordination difficulty (NEXP-complete even cooperatively) and incentive misalignment (publication externalities under individual rewards).
%     \item The publication externality leads to unbounded Price of Anarchy under credit-based rewards without a market.
%     \item The bilateral contract market internalizes the publication externality by enabling agents to trade claims on sequent resolution, creating direct financial incentives for proving and publishing useful results.
%     \item The VAMP game (Decentralized MATP + market) has bounded PoA under sufficient liquidity, but computing exact equilibria is PPAD-hard and the state space is exponentially large.
% \end{enumerate}

These intractability results motivate \emph{learning-based approaches} to approximating equilibria in the VAMP game. Rather than computing exact Nash equilibria (which is provably hard), one may deploy multi-agent reinforcement learning (MARL) algorithms that converge to approximate equilibria through self-play. The VAMP's verifiable settlement mechanism (public resolution of sequents triggers deterministic payoffs) provides a clean reward signal for learning, and the market's price signals serve as a coordination mechanism that reduces the effective dimensionality of the strategy space. We develop this approach in the following chapter, where we instantiate the VAMP framework with specific MARL algorithms (inner-loop MAPPO self-play, outer-loop population-based training approximating PSRO) applied to a synthetic theorem-proving market.